We use the unit-intensity extreme-value normalization: for
\(Z\sim\cN(0,1)\), choose \(u_n\) by \(n\mathbb P\{Z>u_n\}=1\).
Throughout this section, assume \(0\le\eps_n<1\) and write
\[
    X_i=\sigma_n Z_i,\qquad
    \sigma_n=\sqrt{1-\eps_n},\qquad
    Z_i\stackrel{\mathrm{i.i.d.}}{\sim}\cN(0,1).
\]
For a candidate shift \(a\), put \(t=a\sigma_n\). The exact KKT field contains
an empirical-centering factor, which we restore below. Without that factor,
the field is
\[
    e^{-\lambda_n(t/u_n)^2}
    \frac1n\sum_{i=1}^n
    \exp\!\Bigl(tZ_i-\frac{t^2}{2}\Bigr),
    \qquad
    \lambda_n=\frac{\eps_nu_n^2}{2(1-\eps_n)}.
\]
The exponential sum is the correctly specified field; the factor
\(e^{-\lambda_n(t/u_n)^2}\) is the deterministic penalty caused by variance
misspecification. Since $u_n^2/2=\log n+o(\log n)$, $\lambda_n$ and
$\eps_n$ have the same polynomial exponent whenever $\eps_n\to0$. The
one-atom question is whether the penalized field can exceed the KKT level one.

We write \(t=\theta u_n\). After division by the scale
\(c_n(\theta)\) defined below, only observations within distance
\(O(1/u_n)\) of the sample maximum contribute at order one. These
observations need not lie near the score location \(\theta u_n\). Writing an
extreme observation as \(Z_i=u_n+x/u_n\), its contribution is
\[
    \frac1n
    \exp\!\Bigl(\theta u_n^2-\frac{\theta^2u_n^2}{2}\Bigr)e^{\theta x}
    =
    c_n(\theta)e^{\theta x}
\]
to the exponential sum, where
\[
    c_n(\theta)
    :=
    \frac1n
    \exp\!\Bigl(\theta u_n^2-\frac{\theta^2u_n^2}{2}\Bigr)
    =
    n^{-(1-\theta)^2+o(1)}.
\]
For $\theta>1/2$, the centered fluctuation of the correctly specified field is
of order $c_n(\theta)$ and is governed by the Poisson process of extreme
points. Proposition~\ref{prop:poisson-edge-field} proves this convergence
uniformly for $\theta$ in any fixed compact interval
$J\Subset(1/2,1)$.

If \(\lambda_n=n^{-\alpha+o(1)}\), balancing the damping against
\(c_n(\theta)\) gives \(\theta_\alpha=1-\sqrt\alpha\). To pass from compact
convergence to the phase theorem, we must also combine the two sample tails
and exclude crossings below \(1/2\) and near \(1\). The estimates below do so
and establish the needed regularity of the limiting supremum.

\subsection{Compact edge process}

The first input to the polynomial sandwich in
Theorem~\ref{thm:polynomial-stable-regime} is compact convergence of the
normalized empirical exponential sum to the compensated Poisson field.
\begin{proposition}
\label{prop:poisson-edge-field}
Let $Z_1,\dots,Z_n$ be i.i.d.\ $\cN(0,1)$, and let $u_n$ be defined by
\[
    n\mathbb P\{Z_1>u_n\}=1.
\]
For $\theta\in(0,1)$ set
\[
    H_n(\theta)
    =
    \frac1n\sum_{i=1}^n
    \exp\!\Bigl(\theta u_nZ_i-\frac{\theta^2u_n^2}{2}\Bigr),
    \qquad
    c_n(\theta)
    =
    \frac1n
    \exp\!\Bigl(\theta u_n^2-\frac{\theta^2u_n^2}{2}\Bigr).
\]
If $J=[\theta_0,\theta_1]\Subset(1/2,1)$, then
\[
    \left\{
    \frac{H_n(\theta)-1}{c_n(\theta)}
    :
    \theta\in J
    \right\}
    \Rightarrow
    \{S(\theta):\theta\in J\}
    \qquad\text{in }C(J),
\]
where \(N\) is a Poisson point process on \(\R\) with intensity
\(\mu(dx)=e^{-x}\,dx\), and
\[
\begin{aligned}
    S(\theta)&=S^{\rm neg}(\theta)+S^{\rm pos}(\theta),\\
    S^{\rm neg}(\theta)
    &=\lim_{M\to\infty}
      \int_{-M}^0 e^{\theta x}\{N-\mu\}(dx),\\
    S^{\rm pos}(\theta)
    &=\sum_{X\in N\cap[0,\infty)}e^{\theta X}
      -\int_0^\infty e^{\theta x}\mu(dx).
\end{aligned}
\]
The first limit holds in \(L^2(\Omega;C(J))\). The second sum is almost surely
finite because \(\mu([0,\infty))=1\), and its compensator equals
\((1-\theta)^{-1}\). Thus \(S^{\rm pos}(\theta)\) is integrable, though not
square-integrable when \(\theta\ge1/2\).
The process has an almost surely real-analytic version on \((1/2,1)\).
Moreover
\[
    c_n(\theta)
    =
    n^{-(1-\theta)^2+o(1)}
\]
uniformly for $\theta$ in compact subintervals of $(0,1)$.

Let \(X_i=\sigma_n Z_i\), where \(\sigma_n=\sqrt{1-\eps_n}\), and let \(D_n\)
be the centered directional derivative from \eqref{eq:D-intro}. Put
\[
    a_n(\theta)=\frac{\theta u_n}{\sigma_n},
    \qquad
    \lambda_n=\frac{\eps_nu_n^2}{2(1-\eps_n)}.
\]
Then, uniformly on every $J\Subset(1/2,1)$,
\[
    \frac{D_n(a_n(\theta))-1}{c_n(\theta)}
    =
    e^{-\lambda_n\theta^2}
    \frac{H_n(\theta)-1}{c_n(\theta)}
    +
    \frac{e^{-\lambda_n\theta^2}-1}{c_n(\theta)}
    +
    o_{\mathbb P}(1).
\]
\end{proposition}

\begin{proof}
Define the edge point process
\[
    N_n=\sum_{i=1}^n\delta_{u_n(Z_i-u_n)},
    \qquad
    \mu_n=\E N_n.
\]
Its intensity has density
\[
    \mu_n(dx)
    =
    n\frac{\phi(u_n+x/u_n)}{u_n}\,dx
    =
    \frac{n\phi(u_n)}{u_n}
    \exp\!\Bigl(-x-\frac{x^2}{2u_n^2}\Bigr)\,dx.
\]
Since $n\phi(u_n)/u_n\to1$, the standard Poisson convergence theorem for extremes \cite{kallenberg2017random} gives vague convergence of $N_n$ on compact intervals to a Poisson point process $N$ with intensity $e^{-x}dx$. Also
\[
    \frac{H_n(\theta)-1}{c_n(\theta)}
    =
    \int_\R e^{\theta x}\{N_n-\mu_n\}(dx).
\]

Fix \(M<\infty\). Because the limiting intensity is diffuse, the Poisson
process almost surely has no atom at either endpoint of \([-M,M]\). On this
interval, vague convergence, the continuous mapping theorem for the
equicontinuous family $\{e^{\theta x}:\theta\in J\}$, and the uniform
convergence $\mu_n\to\mu$ imply convergence in $C(J)$ of the truncated
centered integrals. We now show that the portions with $x<-M$ and $x>M$
vanish uniformly in \(\theta\) as \(M\to\infty\).

For the upper tail, monotonicity in $\theta$ gives
\[
    \E\sup_{\theta\in J}
    \int_M^\infty e^{\theta x}\{N_n+\mu_n\}(dx)
    \le
    2\int_M^\infty e^{\theta_1x}\mu_n(dx)
    \le
    C_J e^{-(1-\theta_1)M}.
\]
Replacing $N_n,\mu_n$ by their limits $N,\mu$ gives the identical upper-tail
bound.
For the lower tail, write
\[
    R_n(\theta)
    =
    \int_{-\infty}^{-M} e^{\theta x}\{N_n-\mu_n\}(dx).
\]
Since $2\theta_0>1$,
\[
    \E |R_n(\theta_0)|^2
    \le
    \int_{-\infty}^{-M}e^{2\theta_0x}\mu_n(dx)
    \le
    C_J e^{-(2\theta_0-1)M}.
\]
Differentiating in \(\theta\) inserts a factor \(x\); the corresponding
second-moment estimate is
\[
    \sup_{\theta\in J}
    \E |R_n'(\theta)|^2
    \le
    C_J(1+M^2)e^{-(2\theta_0-1)M}.
\]
The squared Sobolev bound
\[
    \|f\|_\infty^2
    \le
    C_J\{|f(\theta_0)|^2+\|f'\|_{L^2(J)}^2\}
\]
on the interval $J$ therefore gives
\[
    \E\sup_{\theta\in J}|R_n(\theta)|^2
    \le
    C_J(1+M^2)e^{-c_JM}\to0
    \qquad(M\to\infty),
\]
uniformly in $n$. The identical estimate with \(N_n,\mu_n\) replaced by
\(N,\mu\) proves the asserted \(L^2(\Omega;C(J))\) construction of
\(S^{\rm neg}\). Together with the upper-tail \(L^1\) bound, this proves the
functional convergence. Repeating the negative-tail estimate after any fixed
number of \(\theta\)-derivatives gives locally uniform convergence of the
derivative series. The positive part is a finite exponential sum minus
\((1-\theta)^{-1}\), so the resulting version of \(S\) is real analytic on
\((1/2,1)\).

Mills' ratio gives
\[
    n
    =
    u_n\sqrt{2\pi}\,e^{u_n^2/2}(1+o(1)).
\]
Hence
\[
    c_n(\theta)
    =
    (u_n\sqrt{2\pi})^{-1}
    \exp\!\Bigl(-\frac{(1-\theta)^2u_n^2}{2}\Bigr)(1+o(1))
    =
    n^{-(1-\theta)^2+o(1)}
\]
uniformly on compact $\theta$-intervals.

It remains to restore the damping and empirical centering in the KKT field.
For $a_n(\theta)=\theta u_n/\sigma_n$,
\[
    D_n(a_n(\theta))
    =
    e^{-\theta u_n\bar Z_n}
    e^{-\lambda_n\theta^2}
    H_n(\theta),
    \qquad
    \bar Z_n=\frac1n\sum_{i=1}^n Z_i.
\]
On $J=[\theta_0,\theta_1]\Subset(1/2,1)$,
\[
    \inf_{\theta\in J}c_n(\theta)
    \ge
    n^{-(1-\theta_0)^2+o(1)}
    \gg n^{-1/2}\sqrt{\log n},
\]
because $(1-\theta_0)^2<1/4$. Since $u_n\bar Z_n=O_{\mathbb P}(\sqrt{\log n/n})$ and
\[
    \sup_{\theta\in J}
    \frac{|H_n(\theta)-1|}{c_n(\theta)}
    =
    O_{\mathbb P}(1),
\]
the empirical-centering factor contributes $o_{\mathbb P}(c_n(\theta))$ uniformly. Dividing by $c_n(\theta)$ gives the stated expansion.
\end{proof}

\subsection{Two tails and empirical centering}

Proposition~\ref{prop:poisson-edge-field} treats positive candidate shifts and
therefore uses the right tail of the sample. Applying it to the reflected
sample $-Z_1,\ldots,-Z_n$ treats negative shifts and the left tail. We now
introduce both signs explicitly.
For $\xi\in\{-1,1\}$ and $\theta\ge0$, set
\[
    a_n(\theta)=\frac{\theta u_n}{\sqrt{1-\eps_n}},
    \qquad
    \lambda_n=\frac{\eps_nu_n^2}{2(1-\eps_n)},
\]
\[
    H_{n,\xi}(\theta)
    =
    \frac1n\sum_{i=1}^n
    \exp\!\Bigl(\theta u_n\xi Z_i-\frac{\theta^2u_n^2}{2}\Bigr),
    \qquad
    \overline D_{n,\xi}(\theta)
    =
    e^{-\lambda_n\theta^2}H_{n,\xi}(\theta),
\]
and
\[
    D_{n,\xi}(\theta)
    =
    D_n\bigl(\xi a_n(\theta)\bigr).
\]
Thus
\[
    D_{n,\xi}(\theta)
    =
    e^{-\xi\theta u_n\bar Z_n}\overline D_{n,\xi}(\theta),
    \qquad
    \bar Z_n=\frac1n\sum_{i=1}^n Z_i.
\]

The exact KKT field $D_{n,\xi}$ differs from $\overline D_{n,\xi}$ only by
the sample-mean factor. Uniformly over both signs and \(0\le\theta\le1\),
this factor changes the field by \(o_{\mathbb P}(\lambda_n)\).

\begin{lemma}
\label{lem:uniform-centering}
Assume $\lambda_n=n^{-\alpha+o(1)}$ for some $\alpha\in(0,1/4)$. Then
\[
    \sup_{\xi=\pm1}\sup_{0\le\theta\le1}
    \bigl|D_{n,\xi}(\theta)-\overline D_{n,\xi}(\theta)\bigr|
    =
    o_{\mathbb P}(\lambda_n).
\]
\end{lemma}

\begin{proof}
From
\[
    D_{n,\xi}(\theta)
    =
    e^{-\xi\theta u_n\bar Z_n}\overline D_{n,\xi}(\theta)
\]
and $u_n|\bar Z_n|=O_{\mathbb P}(\sqrt{\log n/n})=o_{\mathbb P}(1)$, we have
\[
    \sup_{\xi,\theta}
    \bigl|D_{n,\xi}(\theta)-\overline D_{n,\xi}(\theta)\bigr|
    \le
    O_{\mathbb P}(u_n|\bar Z_n|)
    \sup_{\xi,\theta}\overline D_{n,\xi}(\theta).
\]
Thus it is enough to prove
\[
    \sup_{\xi=\pm1}\sup_{0\le\theta\le1}\overline D_{n,\xi}(\theta)
    =
    O_{\mathbb P}(\sqrt{\log n}).
\]
For fixed data, each summand $\exp(\theta u_n\xi Z_i-\theta^2u_n^2/2)$ is maximized over $0\le\theta\le1$ either at $\theta=0$, at $\theta=1$, or at $\theta=\xi Z_i/u_n$ when $0\le\xi Z_i\le u_n$. Hence
\[
    \sup_{0\le\theta\le1}H_{n,\xi}(\theta)
    \le
    1+
    \frac1n\sum_{i=1}^n
    \exp\!\Bigl(\frac{(\xi Z_i)_+^2}{2}\Bigr)
    \mathbf 1_{\{0\le \xi Z_i\le u_n\}}
    +
    H_{n,\xi}(1).
\]
The middle term has expectation $\E[e^{Z^2/2}\mathbf 1_{\{0\le Z\le u_n\}}]=u_n/\sqrt{2\pi}$, so Markov's inequality gives an \(O_{\mathbb P}(u_n)\) bound for the average. The endpoint term $H_{n,\xi}(1)$ is $O_{\mathbb P}(1)$. Indeed, for
\[
    A_n=
    \frac1n\sum_{i=1}^n
    e^{u_n\xi Z_i-u_n^2/2}
    \mathbf 1_{\{\xi Z_i\le u_n\}},
\]
exponential tilting gives $\E A_n=1/2$ and
\[
    \Var(A_n)
    \le
    n^{-1}e^{u_n^2}\Phi(-u_n)
    =
    O(u_n^{-2}).
\]
If \(u_n<\xi Z_i\le u_n+L/u_n\), then the corresponding term in the average, including the factor \(1/n\), is at most \(C_L/u_n\). The number of observations in this window is tight. The probability of any observation above \(u_n+L/u_n\) is \(O(e^{-L})\). Letting \(L\to\infty\) gives \(H_{n,\xi}(1)=O_{\mathbb P}(1)\). Since $e^{-\lambda_n\theta^2}\le1$, the desired $O_{\mathbb P}(\sqrt{\log n})$ bound follows. Multiplying by $u_n|\bar Z_n|$ gives $O_{\mathbb P}(\log n/\sqrt n)=o_{\mathbb P}(\lambda_n)$ because $\alpha<1/4$.
\end{proof}

\subsection{Uniform exclusion and regularity}

Compact convergence does not by itself control the full KKT supremum. The
lemmas in this subsection supply the exclusion and regularity estimates used
to pass from Proposition~\ref{prop:poisson-edge-field} to the polynomial
sandwich in Theorem~\ref{thm:polynomial-stable-regime}. The Poisson
normalization is sharp only for \(\theta>1/2\); below this point,
ordinary fluctuations of the full sample dominate the extreme-tail scale. A
direct second-moment argument controls the larger range
\(\theta\le\theta_*<1/\sqrt2\) at order
\(n^{-1/2+\theta_*^2+o(1)}\). The two methods therefore overlap on
\((1/2,1/\sqrt2)\): the Poisson limit is sharper there, while the cruder
second-moment bound is sufficient for excluding subthreshold crossings.

\begin{lemma}
\label{lem:gaussian-regime-centered}
Fix $0<\theta_*<1/\sqrt2$ and define
\[
    K_{n,\xi}(\theta)
    =
    \frac1n\sum_{i=1}^n
    \exp\!\Bigl(
        \theta u_n\xi(Z_i-\bar Z_n)-\frac{\theta^2u_n^2}{2}
    \Bigr),
    \qquad
    0\le\theta\le\theta_*.
\]
Then
\[
    \sup_{\xi=\pm1}\sup_{0<\theta\le\theta_*}
    \frac{|K_{n,\xi}(\theta)-1|}{\theta^2}
    =
    O_{\mathbb P}\bigl(n^{-1/2+\theta_*^2+o(1)}\bigr).
\]
\end{lemma}

\begin{proof}
We give the proof for $\xi=1$; the other sign is identical. Let
\[
    H_n(\theta)
    =
    \frac1n\sum_{i=1}^n
    \exp\!\Bigl(\theta u_nZ_i-\frac{\theta^2u_n^2}{2}\Bigr).
\]
For $r=1,2,3$, Gaussian integration gives, uniformly for $0\le\theta\le\theta_*$,
\[
    \E\left[
        \left|
        \partial_\theta^r
        \exp\!\Bigl(\theta u_nZ-\frac{\theta^2u_n^2}{2}\Bigr)
        \right|^2
    \right]
    \le
    u_n^{C_r}e^{\theta^2u_n^2}.
\]
Since $\E H_n^{(r)}(\theta)=0$ for $r\ge1$, independence and the displayed
derivative bound imply
\[
    \sup_{0\le\theta\le\theta_*}
    \|H_n^{(r)}(\theta)\|_2
    \le
    u_n^{C_r}n^{-1/2+\theta_*^2+o(1)}.
\]
Applying the elementary Sobolev bound
\[
    \sup_{0\le\theta\le\theta_*}|f(\theta)|
    \le
    |f(0)|+\int_0^{\theta_*}|f'(\theta)|\,d\theta
\]
to $H_n'$ and $H_n''$ gives
\[
    \sup_{0\le\theta\le\theta_*}
    \bigl(|H_n'(\theta)|+|H_n''(\theta)|\bigr)
    =
    O_{\mathbb P}\bigl(n^{-1/2+\theta_*^2+o(1)}\bigr).
\]
Applying the derivative argument once more to $H_n-1$, and using
$H_n(0)=1$, gives
\[
    \sup_{0\le\theta\le\theta_*}|H_n(\theta)-1|
    =
    O_{\mathbb P}\bigl(n^{-1/2+\theta_*^2+o(1)}\bigr).
\]

Write $b_n=u_n\bar Z_n=O_{\mathbb P}(\sqrt{\log n/n})$. Then
\[
    K_{n,1}(\theta)=e^{-b_n\theta}H_n(\theta),
\]
so
\[
    K_{n,1}''(\theta)
    =
    e^{-b_n\theta}
    \{H_n''(\theta)-2b_nH_n'(\theta)+b_n^2H_n(\theta)\}.
\]
The displayed bounds and $b_n=o_{\mathbb P}(1)$ give
\[
    \sup_{0\le\theta\le\theta_*}|K_{n,1}''(\theta)|
    =
    O_{\mathbb P}\bigl(n^{-1/2+\theta_*^2+o(1)}\bigr).
\]
Finally $K_{n,1}(0)=1$ and $K_{n,1}'(0)=0$, hence
\[
    K_{n,1}(\theta)-1
    =
    \int_0^\theta(\theta-s)K_{n,1}''(s)\,ds,
\]
and the asserted bound follows.
\end{proof}

Below \(\theta_\alpha\) by a fixed margin, the deterministic damping
\(\lambda_n\theta^2\) dominates the fluctuations. Gaussian estimates handle
the range up to a fixed $\theta_*>1/2$, and the compact Poisson estimate
handles the remainder.

\begin{lemma}
\label{lem:subthreshold-envelope}
Fix $\alpha\in(0,1/4)$, let
\[
    \theta_\alpha=1-\sqrt\alpha,
\]
and assume $\lambda_n=n^{-\alpha+o(1)}$. For every
\[
    0<\eta<\theta_\alpha-\frac12
\]
we have
\[
    \mathbb P\Bigl\{
        \sup_{\xi=\pm1}
        \sup_{0\le\theta\le\theta_\alpha-\eta}
        D_{n,\xi}(\theta)>1
    \Bigr\}\to0.
\]
\end{lemma}

\begin{proof}
Put $\theta_0=\theta_\alpha-\eta$. Choose
\[
    \theta_* \in
    \Bigl(\frac12,\theta_0\Bigr)
    \quad\text{so that}\quad
    \theta_*^2<\frac12-\alpha;
\]
this is possible because $\theta_0>1/2$ and $\sqrt{1/2-\alpha}>1/2$.

On $[0,\theta_*]$ we use the Gaussian-regime bound for the centered null field. Write
\[
    D_{n,\xi}(\theta)
    =
    e^{-\lambda_n\theta^2}K_{n,\xi}(\theta),
\]
where $K_{n,\xi}$ is the centered null field from Lemma~\ref{lem:gaussian-regime-centered}. Since $\theta_*^2<1/2-\alpha$,
\[
    \sup_{\xi=\pm1}\sup_{0<\theta\le\theta_*}
    \frac{|K_{n,\xi}(\theta)-1|}{\theta^2}
    =
    o_{\mathbb P}(\lambda_n).
\]
Consequently,
\[
    \sup_{\xi=\pm1}\sup_{0<\theta\le\theta_*}
    \frac{D_{n,\xi}(\theta)-1}{\theta^2}
    \le
    -\lambda_n+o_{\mathbb P}(\lambda_n),
\]
because $(1-e^{-\lambda_n\theta^2})/\theta^2=\lambda_n+O(\lambda_n^2)$ uniformly for $0\le\theta\le1$. Therefore $\sup_{\xi=\pm1}\sup_{0\le\theta\le\theta_*}D_{n,\xi}(\theta)\le1$ with probability tending to one.

On the compact interval $[\theta_*,\theta_0]\Subset(1/2,1)$, Lemma~\ref{lem:uniform-centering} lets us replace $D_{n,\xi}$ by $\overline D_{n,\xi}$ at an $o_{\mathbb P}(\lambda_n)$ cost. Proposition~\ref{prop:poisson-edge-field} controls the null edge field uniformly on this compact interval: \(H_{n,\xi}(\theta)-1=O_{\mathbb P}(c_n(\theta))\), uniformly in \(\theta\) and \(\xi\). Moreover,
\[
    \sup_{\theta\le\theta_0}c_n(\theta)
    =
    n^{-(1-\theta_0)^2+o(1)}
    =
    o(\lambda_n),
    \qquad
    1-e^{-\lambda_n\theta^2}
    \ge c_{\theta_*}\lambda_n
\]
for $\theta_*\le\theta\le\theta_0$. We therefore also have
\[
    \sup_{\xi=\pm1}\sup_{\theta_*\le\theta\le\theta_0}
    \overline D_{n,\xi}(\theta)
    \le
    1-c\lambda_n+o_{\mathbb P}(\lambda_n)
    \qquad\text{with probability tending to one}.
\]
Lemma~\ref{lem:uniform-centering} shows that replacing
\(\overline D_{n,\xi}\) by the empirically centered field
\(D_{n,\xi}\) changes the supremum over
\(\xi=\pm1\) and \(\theta_*\le\theta\le\theta_0\) by
\(o_{\mathbb P}(\lambda_n)\). Thus the supremum remains strictly below one
with probability tending to one. \qedhere
\end{proof}

We also need to exclude crossings near \(\theta=1\), where compact
convergence does not apply and the limiting compensator diverges. The next
lemma controls this endpoint uniformly for the finite-sample field. It also
covers the score window of width \(O(u_n^{-2})\) immediately to the right of
\(\theta=1\), beyond the edge score \(t=u_n\).

\begin{lemma}
\label{lem:endpoint-exclusion}
Assume $\lambda_n=n^{-\alpha+o(1)}$ for some $\alpha\in(0,1/4)$. For every fixed $R\in[0,\infty)$,
\[
    \lim_{\rho\downarrow0}\limsup_{n\to\infty}
    \mathbb P\Bigl\{
        \sup_{\xi=\pm1}
        \sup_{1-\rho\le\theta\le1+R/u_n^2}
        D_{n,\xi}(\theta)>1
    \Bigr\}
    =
    0.
\]
\end{lemma}

\begin{proof}
We give the argument for $\xi=1$; the other tail is identical.

We first record the endpoint behavior of the limiting field. Its compensated
negative-half-line part extends continuously to \(\theta=1\). Indeed, for any
\(b>1/2\), the stochastic integrals
\[
 T(\theta)=\int_{-\infty}^0e^{\theta x}\{N-\mu\}(dx),
 \qquad
 T'(\theta)=\int_{-\infty}^0xe^{\theta x}\{N-\mu\}(dx)
\]
have uniformly bounded second moments on \([b,1]\). Stochastic Fubini and the
one-dimensional Sobolev inequality therefore give a continuous version of
\(T\) on \([b,1]\). The positive half-line contains only finitely many
Poisson points, so its random sum has a finite limit as \(\theta\uparrow1\),
whereas its compensator is \((1-\theta)^{-1}\). Consequently,
\begin{equation}
\label{eq:poisson-field-right-endpoint}
 S(\theta)\longrightarrow-\infty
 \qquad\text{almost surely as }\theta\uparrow1.
\end{equation}

We now treat $1-\rho\le\theta<1$, taking \(\rho<1/4\) and
\(\rho^2<\alpha\). The Mills-ratio formula in the proof of
Proposition~\ref{prop:poisson-edge-field} is uniform for
$1-\rho\le\theta\le1$ and gives
\[
    c_n(\theta)
    =
    (u_n\sqrt{2\pi})^{-1}
    \exp\!\left\{
        -\frac{(1-\theta)^2u_n^2}{2}
    \right\}(1+o(1)).
\]
Thus
\[
    \inf_{1-\rho\le\theta<1}c_n(\theta)
    =
    n^{-\rho^2+o(1)}
    \gg \lambda_n.
\]
Since \(c_n(\theta)\ge n^{-\rho^2+o(1)}\) and \(\rho^2<\alpha\),
    \(\lambda_n=o(c_n(\theta))\) uniformly on this interval. Lemma
    \ref{lem:uniform-centering} and the fact that the damping factor is at most
one show that, for every fixed \(\eta>0\), a crossing of \(D_{n,1}\) implies,
with probability tending to one, a crossing above level \(-\eta\) by
\[
    X_n(\theta)
    :=
    \frac{H_{n,1}(\theta)-1}{c_n(\theta)}
    =
    \int_\R e^{\theta x}\{N_n-\mu_n\}(dx).
\]

It remains to control \(X_n\) on a noncompact parameter interval. Split
\([1-\rho,1)\) into
\[
 [1-\rho,1-\rho/2]\quad\text{and}\quad[1-\rho/2,1).
\]
The first interval is compactly contained in \((1/2,1)\), so
Proposition~\ref{prop:poisson-edge-field} and the portmanteau theorem give
\[
 \limsup_{n\to\infty}
 \mathbb P\left\{
   \sup_{1-\rho\le\theta\le1-\rho/2}X_n(\theta)>-\eta
 \right\}
 \le
 \mathbb P\left\{
   \sup_{1-\rho\le\theta\le1-\rho/2}S(\theta)\ge-\eta
 \right\}.
\]
By \eqref{eq:poisson-field-right-endpoint}, the probability on the right
tends to zero as \(\rho\downarrow0\).

For the remaining interval, fix \(M>0\) and write
\[
 A_{n,M}(\theta)
 =
 \int_{-\infty}^M e^{\theta x}\{N_n-\mu_n\}(dx),
 \qquad
 C_{n,M}(\theta)
 =
 \int_M^\infty e^{\theta x}\mu_n(dx).
\]
For \(3/4\le\theta\le1\), the Poisson isometry and the intensity formula in
Proposition~\ref{prop:poisson-edge-field} give
\[
 \sup_{n,\theta}\mathbb E\{|A_{n,M}(\theta)|^2
       +|A_{n,M}'(\theta)|^2\}
 \le C_M.
\]
The Sobolev inequality therefore makes
\(\{\|A_{n,M}\|_\infty:n\ge1\}\) tight. Given
\(\beta>0\), choose \(K=K(M,\beta)\) so that
\[
 \sup_n\mathbb P\{\|A_{n,M}\|_\infty>K\}<\beta.
\]
On \(\{N_n(M,\infty)=0\}\), we have
\(X_n=A_{n,M}-C_{n,M}\). Since \(C_{n,M}\) is increasing in \(\theta\),
\[
 \inf_{1-\rho/2\le\theta<1}C_{n,M}(\theta)
 =
 C_{n,M}(1-\rho/2)
 \longrightarrow
 \frac{2e^{-\rho M/2}}{\rho}
\]
for each fixed \(M,\rho\). This limit exceeds \(K+\eta\) once \(\rho\) is
small enough. Moreover,
\(\limsup_n\mathbb P\{N_n(M,\infty)>0\}\le e^{-M}\). Hence
\[
 \lim_{\rho\downarrow0}\limsup_{n\to\infty}
 \mathbb P\left\{
  \sup_{1-\rho/2\le\theta<1}X_n(\theta)>-\eta
 \right\}
 \le \beta+e^{-M}.
\]
Send first \(\beta\downarrow0\) and then \(M\to\infty\). Together with the
compact interval, this proves the exclusion on \(1-\rho\le\theta<1\).

We now treat the window immediately to the right of one. For each fixed
$R$, we claim that
\[
    \sup_{0\le s\le R}
    H_{n,1}\Bigl(1+\frac{s}{u_n^2}\Bigr)
    \le
    \frac34
    \qquad\text{with probability tending to one}.
\]
Put $t_n(s)=u_n+s/u_n$, and split the average defining $H_{n,1}(1+s/u_n^2)$ at $u_n$. For the bulk term
\[
    B_n(s)
    =
    \frac1n\sum_{i=1}^n
    e^{t_n(s)Z_i-t_n(s)^2/2}
    \mathbf 1_{\{Z_i\le u_n\}},
\]
exponential tilting gives, uniformly for $0\le s\le R$,
\[
    \E B_n(s)
    =
    \mathbb P\{Z\le u_n-t_n(s)\}
    =
    \Phi(-s/u_n)
    \le
    \frac12+o(1),
\]
and
\[
    \Var(B_n(s))
    \le
    \frac1n e^{t_n(s)^2}\Phi(u_n-2t_n(s))
    \le
    \frac{C_R}{u_n^2}.
\]
Differentiating the summand in $s$ gives a factor $(Z_i-t_n(s))/u_n$, and the same tilted calculation yields
\[
    \sup_{0\le s\le R}\Var(B_n'(s))
    \le
    \frac{C_R}{u_n^2}.
\]
Since both \(\sup_s\Var(B_n(s))\) and \(\sup_s\Var(B_n'(s))\) are \(O(u_n^{-2})\), the Sobolev inequality on \([0,R]\) gives \(\sup_{0\le s\le R}|B_n(s)-\E B_n(s)|\to0\) in probability.

For the edge part, first restrict to $u_n<Z_i\le u_n+L/u_n$. On this event each averaged summand is at most $C_{R,L}/u_n$, using $n e^{-u_n^2/2}\asymp u_n$, while the number of such observations is tight; hence this contribution is $o_{\mathbb P}(1)$ for fixed $L$. The probability of any observation above $u_n+L/u_n$ has limit at most $e^{-L}$, so
\[
    \limsup_{n\to\infty}
    \mathbb P\left\{
        \sup_{0\le s\le R}
        \text{edge part}>\varepsilon
    \right\}
    \le
    e^{-L}
\]
for every $\varepsilon>0$. Letting $L\to\infty$ proves the claim.

The damping factor is at most one, and
$e^{-\theta u_n\bar Z_n}=1+o_{\mathbb P}(1)$ uniformly on this window.
Thus the bound for $H_{n,1}$ implies, for example,
$\sup_{1\le\theta\le1+R/u_n^2}D_{n,1}(\theta)\le0.9$ with probability
tending to one. This proves the right-endpoint exclusion. Combining the two
endpoint ranges and then sending $\rho\downarrow0$ completes the proof.
\end{proof}

The phase theorem compares crossing probabilities with events determined by
\(\sup S\). To pass between strict and weak inequalities, we need this
supremum to have no atom at zero. We prove that fact by conditioning on all
but one Poisson point: moving the remaining point changes the supremum
strictly monotonically.

\begin{lemma}
\label{lem:poisson-sup-no-atom}
For every $b\in(1/2,1)$, define
\[
    M_b=\sup_{b\le\theta<1}S(\theta).
\]
Then
\[
    \mathbb P\{M_b=0\}=0.
\]
The same holds for the maximum of two independent copies of $M_b$.
\end{lemma}

\begin{proof}
Use the version of \(S\) that is continuous on compact subintervals of
\([b,1)\) and satisfies \(S(\theta)\to-\infty\) as \(\theta\uparrow1\), as
constructed in the proof of Lemma~\ref{lem:endpoint-exclusion}. For each
bounded interval \(I\), this version can be chosen jointly for every possible
location of a labeled Poisson point in \(I\). Hence \(M_b\) is finite, and its supremum is
attained on a compact subinterval \([b,1-\delta]\).

For $m\ge1$, let
\[
    I_m=[-m-1,-m].
\]
Condition on the Poisson configuration outside $I_m$ and on the number of points inside $I_m$. On the event $N(I_m)=k\ge1$, use the standard labeled representation of the $k$ points in $I_m$ \cite{kallenberg2017random}: after conditioning on $k-1$ labels, the remaining location, call it $Y$, has a continuous density on $I_m$ proportional to $e^{-y}$.

For any conditioned configuration for which the supremum is finite and attained, write the resulting supremum as a function of $Y$:
\[
    \Psi(Y)
    =
    \sup_{b\le\theta<1}
    \{R(\theta)+e^{\theta Y}\},
\]
where $R$ includes all contributions independent of $Y$, including the compensator. If $Y_2>Y_1$ and $\theta_1$ attains the supremum for $Y_1$, then $\theta_1\ge b>0$ and
\[
    \Psi(Y_2)
    \ge
    R(\theta_1)+e^{\theta_1Y_2}
    >
    R(\theta_1)+e^{\theta_1Y_1}
    =
    \Psi(Y_1).
\]
Thus $\Psi$ is strictly increasing, so the equation $\Psi(Y)=0$ has at most one solution. Since $Y$ has a continuous conditional density,
\[
    \mathbb P\{M_b=0\mid N(I_m)=k,\ \text{the conditioned data}\}=0
    \qquad(k\ge1).
\]
After averaging over the conditioning, $\mathbb P\{M_b=0,\ N(I_m)\ge1\}=0$.
Therefore
\[
    \mathbb P\{M_b=0\}
    \le
    \mathbb P\{N(I_m)=0\}
    =
    \exp\{-\mu(I_m)\}.
\]
But $\mu(I_m)=\int_{-m-1}^{-m}e^{-x}\,dx\to\infty$, so $\mathbb P\{M_b=0\}=0$.

If $M_b^+$ and $M_b^-$ are independent copies, then $\mathbb P\{\max(M_b^+,M_b^-)=0\}\le \mathbb P\{M_b^+=0\}+\mathbb P\{M_b^-=0\}=0$.
\end{proof}

It remains to control \(M_b=\sup_{b\le\theta<1}S(\theta)\) as
\(b\downarrow1/2\). The compensated integral over the negative half-line then
accumulates variance of order \((\theta-1/2)^{-1}\). On a sparse geometric
grid, its normalized values approach nearly independent centered Gaussians;
consequently, the probability that every grid value is nonpositive tends to
zero.

\begin{lemma}
\label{lem:poisson-left-endpoint}
For
\[
    M_b=\sup_{b\le\theta<1}S(\theta),
    \qquad
    b\in(1/2,1),
\]
we have
\[
    \mathbb P\{M_b>0\}\to1
    \qquad\text{as }b\downarrow\frac12.
\]
\end{lemma}

\begin{proof}
Split $S(\theta)$ into its negative and positive half-line contributions. On the negative half-line put $r=e^{-x}$. Then the intensity $e^{-x}dx$ becomes Lebesgue measure $dr$ on $[1,\infty)$, and the negative contribution is
\[
    T(\theta)
    =
    \int_1^\infty r^{-\theta}\{N(dr)-dr\},
    \qquad \theta>\frac12,
\]
where $N$ is a unit-rate Poisson process on $[1,\infty)$. The positive half-line contribution is finite and continuous on compact subintervals of $[1/2,1)$, and hence is negligible after multiplication by $\sqrt{\theta-1/2}$. For the finite grid of \(\theta\)-values introduced below, this positive half-line contribution is almost surely bounded and vanishes after multiplication by \(\sqrt{\theta-1/2}\).

We first record the finite-dimensional limit of $T$ near $1/2$. Fix $m\ge1$ and positive constants $a_1,\dots,a_m$. As $\beta\downarrow0$,
\[
    \left\{
        \sqrt{2a_j\beta}\,
        T\Bigl(\frac12+a_j\beta\Bigr)
    \right\}_{j=1}^m
    \Rightarrow
    (G_1,\dots,G_m),
\]
where $(G_1,\dots,G_m)$ is centered Gaussian with covariance
\[
    \E G_iG_j
    =
    \frac{2\sqrt{a_ia_j}}{a_i+a_j}.
\]
Indeed, for any $c_1,\dots,c_m$, write
\[
 g_\beta(r)=
 \sum_{j=1}^m c_j\sqrt{2a_j\beta}\,
 r^{-1/2-a_j\beta}.
\]
Then \(\sup_{r\ge1}|g_\beta(r)|=O(\sqrt\beta)\), while
\[
 \int_1^\infty g_\beta(r)^2\,dr
\]
converges to the quadratic form determined by the displayed covariance.
Expanding the compensated Poisson characteristic exponent gives
\[
 \int_1^\infty
 \{e^{it g_\beta(r)}-1-itg_\beta(r)\}\,dr
 =
 -\frac{t^2}{2}\int_1^\infty g_\beta(r)^2\,dr+o(1),
\]
because the remainder is bounded by
\(C_t\|g_\beta\|_\infty\int g_\beta^2=o(1)\). This proves the stated
Gaussian finite-dimensional limit.

Fix $m$ and $q\in(0,1)$. For $b>1/2$ set
\[
    \beta_b=b-\frac12,
    \qquad
    \theta_j(b)=\frac12+\beta_b q^{-j},
    \qquad
    0\le j\le m-1.
\]
For $b$ close enough to $1/2$, all these points lie in $[b,1)$. Applying the
Gaussian finite-dimensional limit just proved with $a_j=q^{-j}$ gives
\[
    \left\{
        \sqrt{2\beta_b q^{-j}}\,
        S(\theta_j(b))
    \right\}_{j=0}^{m-1}
    \Rightarrow
    (G_0,\dots,G_{m-1}),
\]
where
\[
    \E G_iG_j
    =
    \frac{2q^{|i-j|/2}}{1+q^{|i-j|}}.
\]
Thus
\[
    \limsup_{b\downarrow1/2}\mathbb P\{M_b\le0\}
    \le
    \mathbb P\{G_0\le0,\dots,G_{m-1}\le0\}.
\]
For fixed $m$, the covariance matrix above converges to the identity as $q\downarrow0$. Hence the right-hand side tends to $2^{-m}$ as $q\downarrow0$. Letting first $q\downarrow0$ and then $m\to\infty$ proves the lemma.
\end{proof}
